% Zone „Das kennst du schon“ – aus bank/quadratische-gleichungen/zone.jsonl (zusammenbau v0.1)
\einheitenkopf[zone][Kennst du schon]{Das kennst du schon}
\setcounter{aufgabe}{0}
%% TODO Grundvorstellungs-Aufgabe als erste Hauptnummer der Zone (2.8) fehlt in der Bank

%% TODO Titel: Ich-kann-Satz fehlt in der Bank (Platzhalter: Kettenname) – Nr. 1
%% TODO Anweisung: ein Satz über der ersten Teilaufgabe fehlt in der Bank – Nr. 1
\begin{aufgabe}{Quadrieren auch negativer Zahlen und Dezimalzahlen, Quadrat vor Punkt vor Strich}
%% TODO Darstellung für schwach fehlt in der Bank (quadratische-gleichungen-zone-f1-v1; Abschnitt „Für schwache Schüler“)
\swz{$7^2$ – wie viel?}{}
%% TODO Darstellung für schwach fehlt in der Bank (quadratische-gleichungen-zone-f1-v3; Abschnitt „Für schwache Schüler“)
\swz{$(-1{,}2)^2$ – wie viel?}{}
\end{aufgabe}

%% TODO Titel: Ich-kann-Satz fehlt in der Bank (Platzhalter: Kettenname) – Nr. 2
%% TODO Anweisung: ein Satz über der ersten Teilaufgabe fehlt in der Bank – Nr. 2
\begin{aufgabe}{Quadratwurzel ziehen}
%% TODO Darstellung für schwach fehlt in der Bank (quadratische-gleichungen-zone-f2-v1; Abschnitt „Für schwache Schüler“)
\swz{$\sqrt{49}$ – wie viel?}{}
%% TODO Darstellung für schwach fehlt in der Bank (quadratische-gleichungen-zone-f2-v3; Abschnitt „Für schwache Schüler“)
\swz{$\sqrt{11}$ – auf zwei Stellen?}{}
\end{aufgabe}

%% TODO Titel: Ich-kann-Satz fehlt in der Bank (Platzhalter: Kettenname) – Nr. 3
%% TODO Anweisung: ein Satz über der ersten Teilaufgabe fehlt in der Bank – Nr. 3
\begin{aufgabe}{Lineare Gleichung mit Äquivalenzumformungen und Strich lösen, auch negative Vorzahl und x auf beiden Seiten}
%% TODO Darstellung für schwach fehlt in der Bank (quadratische-gleichungen-zone-f3-v1; Abschnitt „Für schwache Schüler“)
\swz{$\text{x + 8 = 15}$}{}
%% TODO Darstellung für schwach fehlt in der Bank (quadratische-gleichungen-zone-f3-v3; Abschnitt „Für schwache Schüler“)
\swz{$\text{-2x = 14}$}{}
\end{aufgabe}

%% TODO Titel: Ich-kann-Satz fehlt in der Bank (Platzhalter: Kettenname) – Nr. 4
%% TODO Anweisung: ein Satz über der ersten Teilaufgabe fehlt in der Bank – Nr. 4
\begin{aufgabe}{Lösung durch Einsetzen prüfen mit (wA)/(fA), auch mit negativer Zahl}
%% TODO Darstellung für schwach fehlt in der Bank (quadratische-gleichungen-zone-f4-v1; Abschnitt „Für schwache Schüler“)
\swz{Ist $x = 4$ Lösung von $x + 9 = 13$?}{}
%% TODO Darstellung für schwach fehlt in der Bank (quadratische-gleichungen-zone-f4-v3; Abschnitt „Für schwache Schüler“)
\swz{Ist $x = -6$ Lösung von $x + 10 = 4$?}{}
\end{aufgabe}

%% TODO Titel: Ich-kann-Satz fehlt in der Bank (Platzhalter: Kettenname) – Nr. 5
%% TODO Anweisung: ein Satz über der ersten Teilaufgabe fehlt in der Bank – Nr. 5
\begin{aufgabe}{Klammer zuerst und Punkt vor Strich mit negativen Zahlen}
%% TODO Darstellung für schwach fehlt in der Bank (quadratische-gleichungen-zone-f5-v1; Abschnitt „Für schwache Schüler“)
\swz{$3 \cdot (2 + 4)$ – wie viel?}{}
%% TODO Darstellung für schwach fehlt in der Bank (quadratische-gleichungen-zone-f5-v3; Abschnitt „Für schwache Schüler“)
\swz{$(-4) \cdot (-4 + 9)$ – wie viel?}{}
\end{aufgabe}

%% TODO Titel: Ich-kann-Satz fehlt in der Bank (Platzhalter: Kettenname) – Nr. 6
%% TODO Anweisung: ein Satz über der ersten Teilaufgabe fehlt in der Bank – Nr. 6
\begin{aufgabe}{Scheitelpunktform lesen}
%% TODO Darstellung für schwach fehlt in der Bank (quadratische-gleichungen-zone-f6-v1; Abschnitt „Für schwache Schüler“)
\swz{$y = (x - 2)^2 + 1$ – Scheitel?}{}
%% TODO Darstellung für schwach fehlt in der Bank (quadratische-gleichungen-zone-f6-v3; Abschnitt „Für schwache Schüler“)
\swz{$y = (x + 1)^2 - 6$ – Scheitel?}{}
\end{aufgabe}

%% TODO Titel: Ich-kann-Satz fehlt in der Bank (Platzhalter: Kettenname) – Nr. 7
%% TODO Anweisung: ein Satz über der ersten Teilaufgabe fehlt in der Bank – Nr. 7
\begin{aufgabe}{Ausklammern, Ausmultiplizieren und binomische Formeln}
%% TODO Darstellung für schwach fehlt in der Bank (quadratische-gleichungen-zone-f7-v1; Abschnitt „Für schwache Schüler“)
\swz{$x \cdot (x + 4)$ – ausmultiplizieren?}{}
%% TODO Darstellung für schwach fehlt in der Bank (quadratische-gleichungen-zone-f7-v3; Abschnitt „Für schwache Schüler“)
\swz{$(x - 6)^2$ – ausmultiplizieren?}{}
\end{aufgabe}

%% TODO Titel: Ich-kann-Satz fehlt in der Bank (Platzhalter: Kettenname) – Nr. 8
%% TODO Anweisung: ein Satz über der ersten Teilaufgabe fehlt in der Bank – Nr. 8
\begin{aufgabe}{Terme ordnen und zusammenfassen, x² zuerst, dann x-Glieder, dann Zahlen}
%% TODO Darstellung für schwach fehlt in der Bank (quadratische-gleichungen-zone-f8-v1; Abschnitt „Für schwache Schüler“)
\swz{$3x + x^2 + 2$ – ordnen?}{}
%% TODO Darstellung für schwach fehlt in der Bank (quadratische-gleichungen-zone-f8-v3; Abschnitt „Für schwache Schüler“)
\swz{$x^2 - 8 + 3x + 2$ – ordnen und zusammenfassen?}{}
\end{aufgabe}

%% TODO Titel: Ich-kann-Satz fehlt in der Bank (Platzhalter: Kettenname) – Nr. 9
%% TODO Anweisung: ein Satz über der ersten Teilaufgabe fehlt in der Bank – Nr. 9
\begin{aufgabe}{Quadrieren auch negativer Zahlen und Dezimalzahlen, Quadrat vor Punkt vor Strich – Fehler finden}
\swz[3]{Wo ist der Fehler? $(-8)^2 = -64$}{}
\end{aufgabe}

%% TODO Titel: Ich-kann-Satz fehlt in der Bank (Platzhalter: Kettenname) – Nr. 10
%% TODO Anweisung: ein Satz über der ersten Teilaufgabe fehlt in der Bank – Nr. 10
\begin{aufgabe}{Quadrieren auch negativer Zahlen und Dezimalzahlen, Quadrat vor Punkt vor Strich – selbst rechnen}
%% TODO Darstellung für schwach fehlt in der Bank (quadratische-gleichungen-zone-f1-v6; Abschnitt „Für schwache Schüler“)
\swz{$(-13)^2$ – wie viel?}{}
\end{aufgabe}
